% D.79. Root, FGH, and PBC complete independent proof/source/certificate audits PASS.
\subsection{The entire compact kernel on the two-prime invertible matrix family}

\begin{statement}{D.79}{Proved}{The complete two-prime matrix-family compact kernel}
Let \(M\in\operatorname{GL}_2(\mathbb Z_5)\), and put
\[
 \binom ab=M\binom{\kappa(2)}{\kappa(3)},\qquad
 \binom\alpha\beta
 =M\binom{\log16/(4\log6)}{\log81/(4\log6)}.
                                                               \tag{D.79.1}
\]
The vector \((\alpha,\beta)\) is primitive. With
\(v_5(0)=\infty\), set
\[
 d=\max\{v_5(\alpha),v_5(\beta),v_5(\alpha-\beta)\}.
                                                               \tag{D.79.2}
\]
At most one of these valuations is positive. For the actual
Alexander twist \(s=(a,b)\), we have
\[
 \mathcal S_s=
 \begin{cases}
  \mathbb Z/5,&d=0,\\
  0,&1\le d\le\infty,
 \end{cases}
 \qquad \mathcal T_s=0,
                                                               \tag{D.79.3}
\]
and, as compact groups,
\[
 \mathcal K_s\simeq
 \begin{cases}
  \mathbb Z/5,&d=0,\\
  \mathbb Z_5^d,&1\le d<\infty,\\
  \displaystyle\prod_{\mathbb N}\mathbb Z_5,&d=\infty.
 \end{cases}                                                   \tag{D.79.4}
\]
Here \(\mathcal S_s\) is exactly the inverse finite-global-localization
kernel of D.63, not a kernel formed only after passing to divisible
local coefficients. The separated second cohomology is the open
extension of the actual residue space by (D.79.4) in D.63.
No splitting of that latter extension is asserted.

This classifies the compact kernel for every matrix in this family,
including arbitrarily deep and exact boundary slopes. It does not
assert that every matrix pair has a metabelian section. If its fibre
is nonempty, D.74 makes it a torsor under the compact countable
product of \(\mathbb Z_5\). The identity matrix includes the actual
point \(t=-2\) of D.78.
\end{statement}

\begin{proof}
Use D.78's notation \(I_0,I_s,H,D,\lambda\). Inertia at two and
three gives the two columns of \(M\) in the translation group.
Because \(M\) is invertible, these columns generate the whole
\(H=\mathbb Z_5^2\). The faithful actual operator argument of
D.78 therefore gives
\[
 I_s=H\rtimes I_0.
                                                               \tag{D.79.5}
\]
The local surjection \(G_{\mathbb Q_5}\to I_0\) is independent
of the pair. After the same integral Kummer gauge as in D.78,
the local pair is \((\alpha\kappa(6),\beta\kappa(6))\).
The vector is primitive since the vector \((2,4)\bmod5\) is
nonzero and \(M\) is invertible. D.78's exact calculation
\(\kappa(6)(\ker(G_{\mathbb Q_5}\to I_0))=\mathbb Z_5\)
now gives the actual equality
\[
 I_5=H_5\rtimes I_0,
 \qquad H_5=\mathbb Z_5(\alpha,\beta).
                                                               \tag{D.79.6}
\]
Thus continuous translation coinduction gives, naturally,
\[
 H^1(I_s,D)=H^1(I_0,\lambda),\qquad
 H^1(I_5,D)=H^1(I_0,D_{\rm lin}),
 \quad\lambda=\mathbb Q_5/\mathbb Z_5(-1).
                                                               \tag{D.79.7}
\]
Every finite everywhere locally trivial class descends to the actual
operator image by D.78's finite affine Chebotarev argument.
Consequently
\[
 \mathcal S_s^*\ \subseteq\
 \ker\{H^1(I_0,\lambda)\longrightarrow H^1(I_0,D_{\rm lin})\}.
                                                               \tag{D.79.8}
\]
Here the star is Pontryagin dual; all coefficient modules inside
the cohomology groups are Cartier duals.

\paragraph{Boundary permutations are actual geometric comparisons.}
The rational automorphisms \(t\mapsto1-t\) and \(t\mapsto1/t\)
of the tripod permute its three peripheral directions. Their
maps on pro-five Kummer pairs are
\[
 (a,b)\longmapsto(b,a),\qquad
 (a,b)\longmapsto(-a,b-a),                              \tag{D.79.9}
\]
because the extra Kummer class of \(-1\) is zero at five.
These are comparisons of the actual arithmetic fundamental
extensions. Choose a geometric path between the transported unit
tangent and the original reference tangent. Its ambiguity is an
inner geometric automorphism, which gives a coefficient-module
gauge and hence the same continuous cohomology groups. The
transported unit cusp has zero abelian class, by the three actual
packet formulas in D.35; an integral abelian gauge therefore
identifies the reference abelian representatives. Thus (D.79.9)
transports the actual Alexander module and every local and global
cohomology map, including its finite-coefficient inverse limits.
It is not an inference from an abstract formal substitution.

The transformations are invertible over \(\mathbb Z_5\), keep
ramification inside \(\{2,3,5\}\), and permute the three
valuations in (D.79.2). They allow us, whenever \(d>0\), to
arrange that
\[
 v_5(\alpha)=d,\qquad \beta\in\mathbb Z_5^\times.
                                                               \tag{D.79.10}
\]
The exact diagonal comparison could equivalently be made using
D.39's half-cyclotomic monomial gauge, since two is a unit.

\paragraph{The case of three unit slopes.}
If \(d=0\), the proof of D.78 applies verbatim to the actual
images (D.79.5)--(D.79.6): the primitive local cubic coefficient
is \(\kappa_3\beta(-\alpha)(\beta-\alpha)/48\), a unit.
It computes the entire restriction kernel in (D.79.8) as
\(\mathbb F_5\). Its generator is the reduction of the integral
cocycle \(48h\) of D.31. The fixed ten-dimensional module of
D.78 still trivializes it at every prime. Outside two, three,
five everything is unramified; at two and three the constant
primitive works because the residual cyclotomic multipliers
are not one. The local unit-slope three-jet proof supplies
the primitive at five. Hence this is a single finite-level
everywhere locally trivial class, and
\(\mathcal S_s=\mathbb Z/5\).

\paragraph{An exact boundary gives an injective local restriction.}
If \(d=\infty\), make (D.79.10) with \(\alpha=0\).
Then
\[
 D_{\rm lin}=(\Lambda/(t_2-1))^D.
\]
Its action factors through \(\Gamma=\mathbb Z_5^\times\),
by the exact axis formula \(\chi S_\chi(t_1)\) in D.35.
Put \(N_0=\ker(I_0\to\Gamma)\). Both \(\lambda\) and
\(D_{\rm lin}\) have trivial \(N_0\)-action. The inclusion
\(\lambda\subset D_{\rm lin}\) is injective, so it induces
an injection on continuous homomorphisms from \(N_0\).
Moreover
\(H^1(\Gamma,\lambda)=0\), since its order-four subgroup
acts by a nontrivial scalar with unit difference from one.
Inflation--restriction therefore makes
\(H^1(I_0,\lambda)\to H^1(I_0,D_{\rm lin})\) injective.
Equation (D.79.8) gives \(\mathcal S_s=0\).

\paragraph{The exact finite module at a boundary depth.}
Now assume \(1\le d<\infty\) and (D.79.10). Write
\[
 N=5^d,\qquad \alpha=N a_0,
 \quad a_0,\beta\in\mathbb Z_5^\times,
 \qquad R=\mathbb Z_5[[z]],\quad t=1+z.
\]
The local line quotient is given integrally by
\(t_1=t^\beta,t_2=t^{-\alpha}\). Define the distinguished
polynomial
\[
 Q(z)=z^2(t^N-1),\qquad \deg Q=N+2.                   \tag{D.79.11}
\]
For every actual cyclotomic-kernel element \(\tau\in I_0\),
D.39's three-boundary factorization shows that its multiplier
minus one is divisible by \(Q\). For the actual primitive
local \(\delta\) of D.77 it is exactly
\[
                   \rho_0(\delta)-1=Q(z)E(z),
                          \qquad E\in R^\times.        \tag{D.79.12}
\]
Indeed \(t^{-\alpha}-1\) is \(t^N-1\) times a unit,
whereas \(t^\beta-1\) and \(t^{\beta-\alpha}-1\)
are \(z\) times units. The remaining beta factor is a unit
by the primitive cubic coefficient. All these statements are
integral, with no division by \(5^d\).

Consequently
\[
 D_{\rm lin}^{I_0}=(R/QR)^D{}^{\Gamma},\qquad
 (D_{\rm lin}/\lambda)^{I_0}
                    =(zR/zQR)^D{}^{\Gamma}.           \tag{D.79.13}
\]
The action on both finite-rank compact quotients factors
through \(\Gamma\), because every \(\tau\in N_0\)
acts by a multiplier congruent to one modulo \(Q\).
The second identity uses the dual of the stable ideal \(zR\).

Let \(V=(zR/zQR)^D\), and put \(J=(R/QR)^D\).
The actual inclusion of ideals gives the exact dual sequence
\[
 0\longrightarrow\lambda\longrightarrow J
       \longrightarrow V
       \longrightarrow (QR/zQR)^D\longrightarrow0.
                                                               \tag{D.79.14}
\]
On \(QR/zQR\), the arithmetic character is \(\chi^5\):
the multiplier has constant term \(\chi^2\), and the
ratio \(Q(t^\chi-1)/Q(t-1)\) has constant term \(\chi^3\).
The last Cartier dual is therefore
\(\mathbb Q_5/\mathbb Z_5(-4)\). Since
\(H^0(\Gamma,\lambda)=H^1(\Gamma,\lambda)=0\),
(D.79.14) yields an injection
\[
 \frac{V^\Gamma}{D_{\rm lin}^{I_0}}
       \lhook\joinrel\longrightarrow
 (\mathbb Q_5/\mathbb Z_5(-4))^\Gamma
                    =\mathbb F_5.                    \tag{D.79.15}
\]
The quotient on the left is exactly the full restriction
kernel in (D.79.8), by the coefficient long exact sequence.
Thus it remains to exclude this one possible order-five class.

\paragraph{A uniform actual cyclotomic operator kills that class.}
Let \(k\in I_0\) be actual with \(\chi(k)=2\). For a general
actual element with cyclotomic multiplier \(\gamma\), write its
operator on the line as \(U_\gamma(t)\varphi_\gamma\), where
\(\varphi_\gamma(t)=t^\gamma\). Its class modulo \(Q\) is
independent of the choice of that element, by (D.79.12).
The exact axis formula shows that
\(U_\gamma-\gamma S_\gamma(t^\beta)\) is divisible by
\(t^N-1\): each of its two differences contains
\(t^{-\alpha}-1\). Modulo \(Q=z^2(t^N-1)\) it therefore is
\[
 U_\gamma(t)\equiv
 \gamma S_\gamma(t^\beta)+(t^N-1)(A_\gamma+B_\gamma z)
                                                      \pmod Q,\tag{D.79.16}
\]
where the exact first two Taylor coefficients give
\[
 \begin{split}
 A_\gamma&=-\frac{a_0\gamma^2(\gamma-1)}2,\\
 B_\gamma&=a_0\gamma^2(\gamma-1)
       \left(\frac{\alpha(2\gamma-1)}{12}
                    +\frac N4-\frac{\beta(7\gamma-5)}{24}\right).
 \end{split}                                             \tag{D.79.17}
\]
To check the only beta input in this computation, its degree-two
coefficient is \((\gamma^2-1)uv/24\), exactly D.31/D.39.
All denominators in (D.79.17) are five-adic units. In particular
\(A_\gamma,B_\gamma\in5\mathbb Z_5\) whenever
\(\gamma\equiv1\pmod5\).

Take the actual power
\[
 g=k^{4\cdot5^{d-1}},\qquad
 \gamma=2^{4\cdot5^{d-1}}
                  \equiv1+3N\pmod{5N}.               \tag{D.79.18}
\]
The congruence follows by binomial expansion of
\((1+15)^{5^{d-1}}\). Modulo five, to degree \(N+2\),
\[
 \begin{split}
 \varphi_\gamma(z)&\equiv z+3z^N+3z^{N+1},\\
 U_\gamma(t)&\equiv S_\gamma(t^\beta)
                 \equiv1+3z^{N-1}+O(z^N)\pmod Q.
 \end{split}                                             \tag{D.79.19}
\]
For the second line, use \(\beta^{N-1}\equiv1\pmod5\).
The first line follows from
\((1+z)^\gamma\equiv(1+z)(1+3z^N)\) in this range;
the omitted next degree \(2N\) exceeds \(N+2\), since
\(N\ge5\).

In the compact module \(zR/zQR\), whose reduction modulo five
has basis \(z,\ldots,z^{N+2}\), these identities imply
\[
 (\gamma^{-1}\rho_0(g)-1)z^3
       \equiv (3\cdot3+3)z^{N+2}
       =2z^{N+2}\pmod{5,zQ}.                           \tag{D.79.20}
\]
The leading term of \(Q\) modulo five is exactly \(z^{N+2}\).
Every order-five \(\Gamma\)-invariant Cartier functional on
\(zR/zQR\) annihilates the left side of (D.79.20); hence it
annihilates \(Q\). Its image in the last term of (D.79.14)
is therefore zero. This argument uses an actual cyclotomic
power, not a formal graded eigenvalue.

Finally D.77 computes
\[
 D_{\rm lin}^{I_0}=D^{I_5}
                  \simeq(\mathbb Q_5/\mathbb Z_5)^d.   \tag{D.79.21}
\]
It is divisible. If the quotient in (D.79.15) were nonzero,
the inclusion of this divisible group into \(V^\Gamma\)
would split abstractly: a divisible abelian group is injective.
Its quotient, which has order five, would consequently have
an order-five lift in \(V^\Gamma\) mapping nontrivially to
\((QR/zQR)^D\). Equation (D.79.20) excludes that lift.
Thus the whole restriction kernel in (D.79.8) is zero.
This proves \(\mathcal S_s=0\) for every finite positive depth.

\paragraph{The compact groups and the remaining scope.}
Every one of \(a,b,a-b\) has global divisibility zero: the two
rows of an invertible matrix, and their difference, are
primitive. The local depths are respectively the valuations
in (D.79.2). D.77 therefore gives \(\mathcal T_s=0\),
\(L_5=\mathbb Z_5^d\) at finite depth, and the countable
compact product at infinite depth. Its continuous product
decomposition \(\mathcal K_s=\mathcal S_s\times L_5\)
proves (D.79.4). Finiteness of the ramification set is automatic
for the completed matrix combinations; D.74 applies to any
nonempty fibre. Section existence itself and the separated
extension class are not deduced from the cohomology calculation.
\end{proof}
